\documentclass[12pt,a4paper,oneside]{article}

\usepackage[utf8]{inputenc}
\usepackage[T1]{fontenc}
\usepackage[brazil]{babel}
\usepackage{mathptmx}
\usepackage[a4paper,top=3cm,bottom=3cm,left=3cm,right=3cm]{geometry}
\usepackage{setspace}
\onehalfspacing
\usepackage{indentfirst}
\setlength{\parindent}{1.25cm}
\usepackage{graphicx}
\usepackage{enumitem}
\usepackage{hyperref}
\hypersetup{hidelinks}

\begin{document}

% =============================================================
% CAPA
% =============================================================
\begin{center}
{\large Universidade Federal de Itajubá -- Campus de Itabira\\[0.4cm]
Instituto de Ciências e Tecnologias}\\[3cm]

{\Huge\bfseries RELATÓRIO FINAL DE ESTÁGIO}\\[3cm]

\begin{flushleft}
Discente: Pedro Henrique Leite Mendes\\
RA: 2020030824\\
Curso: Engenharia de Controle e Automação\\
Data de Entrega: 25/05/2026
\end{flushleft}
\end{center}

\thispagestyle{empty}
\clearpage

% =============================================================
% SUMÁRIO
% =============================================================
\renewcommand{\contentsname}{Conteúdo}
\tableofcontents
\thispagestyle{empty}
\clearpage

\setcounter{page}{1}

% =============================================================
% 1. INTRODUÇÃO
% =============================================================
\section{Introdução}

\subsection{Objetivo}

O objetivo deste estágio em Engenharia de Controle e Automação na AwSales LTDA foi adquirir experiência prática e aplicar os conhecimentos teóricos adquiridos durante o curso de graduação em situações reais de trabalho, atuando na configuração, monitoramento e otimização de agentes de Inteligência Artificial aplicados a fluxos de comunicação e atendimento digital.

\subsection{Dados do aluno estagiário}

\begin{itemize}[leftmargin=2em,topsep=0pt]
  \item Nome: Pedro Henrique Leite Mendes.
  \item Curso: Engenharia de Controle e Automação.
  \item Instituição de Ensino: Universidade Federal de Itajubá -- Campus Itabira.
  \item Matrícula: 2020030824.
  \item Contato: pedrohlmendes@hotmail.com.
\end{itemize}

\subsection{Dados da Empresa onde o estágio foi realizado}

\begin{itemize}[leftmargin=2em,topsep=0pt]
  \item Nome da Empresa: AwSales LTDA.
  \item Setor de Atuação: Desenvolvimento e licenciamento de programas de computador (\textit{Autonomous Service Builder} de agentes de Inteligência Artificial).
  \item Endereço: Rua Ministro Orozimbo Nonato, 102 -- Vila da Serra, Nova Lima -- MG, 34006-053.
  \item CNPJ: 59.298.768/0001-36.
  \item Site: \url{https://www.awsales.io/}.
\end{itemize}

\subsection{Período de realização do estágio}

\begin{itemize}[leftmargin=2em,topsep=0pt]
  \item Data de Início: 01/10/2025.
  \item Data de Término: 25/05/2026.
  \item Carga horária: 30 horas semanais, totalizando aproximadamente 966 horas (novecentas e sessenta e seis horas).
\end{itemize}

\subsection{Dados do Supervisor do Estágio}

\begin{itemize}[leftmargin=2em,topsep=0pt]
  \item Nome do Supervisor: Daniel Roscoe Oliveira.
  \item Cargo: Gerente de Estratégia e Operações.
  \item Contato: daniel@mambaculture.com.
\end{itemize}

% =============================================================
% 2. ATIVIDADES DESENVOLVIDAS
% =============================================================
\section{Atividades desenvolvidas}

\subsection{Descrição das atividades desenvolvidas pela empresa}

A AwSales LTDA é uma empresa brasileira de tecnologia sediada em Nova Lima -- MG, que se posiciona como um \textit{Autonomous Service Builder}: uma plataforma de Inteligência Artificial que opera agentes autônomos especializados em vendas, atendimento, retenção e \textit{onboarding}, atuando 24 horas por dia em canais conversacionais como o WhatsApp. A empresa é \textit{Meta Business Partner}, com acesso prioritário às APIs oficiais do WhatsApp, e atende clientes \textit{enterprise} em setores como educação (Uniasselvi, G4 Educação) e segurança residencial (Emive), além de operações digitais de infoprodutores e \textit{e-commerces}.

A plataforma é composta por quatro grandes módulos: \textit{Agent Studio}, que traduz objetivos descritos em linguagem natural em comportamento operacional de agente sem necessidade de fluxogramas; \textit{AgentsCore}, conjunto de arquiteturas estratégicas que definem como o agente pensa, decide e executa; \textit{Memory Base}, base de conhecimento em camadas (contexto institucional, regras operacionais, dados por usuário, histórico de conversa); e \textit{Cortex}, sistema de otimização contínua que monitora conversas em tempo real, testa variações de abordagem em paralelo e promove automaticamente as que geram melhor resultado. A plataforma oferece mais de 100 integrações com CRMs, ERPs, gateways de pagamento e ferramentas de suporte, com promessa de implantação em até 14 dias úteis.

Durante o estágio, o estagiário atuou na área de operação e otimização desses agentes de IA, com foco na parametrização e manutenção de \textit{checkpoints} (instruções de comportamento), em FAQs (base de conhecimento), em \textit{tools} (integrações com sistemas externos), em mensagens de disparo para WhatsApp e na análise de dados de desempenho e custo das campanhas em produção.

\subsection{Cronologia de desenvolvimento das atividades referentes ao estágio}

O desenvolvimento do estágio seguiu a progressão estipulada no Plano de Atividades acordado entre a instituição de ensino e a empresa concedente, sendo estruturado em quatro grandes etapas, distribuídas ao longo de aproximadamente oito meses.

\textbf{Etapa 1 -- Outubro a Novembro de 2025 (Integração e Treinamento):}
\begin{itemize}[leftmargin=2em,topsep=0pt]
  \item Integração na empresa e reuniões iniciais com o supervisor.
  \item Treinamento técnico na plataforma AwSales e em ambiente de testes.
  \item Estudo da documentação interna e da arquitetura multiagente do produto.
  \item Suporte e acompanhamento de clientes ativos, sob supervisão direta.
\end{itemize}

\textbf{Etapa 2 -- Dezembro de 2025 a Fevereiro de 2026 (Criação de Agentes):}
\begin{itemize}[leftmargin=2em,topsep=0pt]
  \item Criação e configuração de agentes de IA em campanhas novas.
  \item Aplicação do fluxo obrigatório de três fases para criação de campanhas (textos complementares, FAQs, checkpoint, mensagens de disparo).
  \item Atendimento a demandas de \textit{Customer Success} e \textit{Account Management}.
  \item Reuniões de acompanhamento com o supervisor.
\end{itemize}

\textbf{Etapa 3 -- Março a Abril de 2026 (Análise e Otimização):}
\begin{itemize}[leftmargin=2em,topsep=0pt]
  \item Análise de dados de desempenho e custos de campanhas em produção.
  \item Otimização de \textit{checkpoints} e redução do consumo de \textit{tokens} dos sub-agentes mais caros.
  \item Verificação cruzada de FAQs para evitar conflitos com as regras do \textit{checkpoint}.
  \item Validação automatizada de formato de \textit{tools} via expressões regulares.
\end{itemize}

\textbf{Etapa 4 -- Maio de 2026 (Integrações e Documentação):}
\begin{itemize}[leftmargin=2em,topsep=0pt]
  \item Refatoração de integrações entre a plataforma AwSales e sistemas externos via \textit{webhooks} HTTPS.
  \item Padronização do contrato de resposta das \textit{tools} críticas (padrão \textit{payload-ok}).
  \item Consolidação da documentação interna em arquivos Markdown versionados em Git.
  \item Elaboração do relatório final e dos documentos de entrega à Coordenação de Estágio.
\end{itemize}

% =============================================================
% 3. METODOLOGIAS E TECNOLOGIAS
% =============================================================
\section{Metodologia e tecnologias envolvidas no desenvolvimento do estágio}

As atividades exigiram domínio da plataforma proprietária AwSales, conhecimento sobre modelos de linguagem (LLMs) utilizados nos sub-agentes da cadeia, ferramentas de orquestração externa (n8n, Zapier, ActiveCampaign), linguagens auxiliares como Python e ferramentas de versionamento (Git). Os tópicos abaixo descrevem as principais frentes técnicas do estágio.

\subsection{Configuração de agentes de IA}

A configuração de um agente na plataforma AwSales é composta por três artefatos principais: \textit{Checkpoint} (instruções de comportamento e fluxo), \textit{Base de Conhecimento} (FAQs de Produto e Playbook consultadas por busca semântica) e \textit{Mensagens de Disparo} (abertura e \textit{follow-ups}). O estagiário acompanhou e executou o ciclo completo de criação de campanhas, seguindo o padrão interno de geração de textos complementares, avaliação das FAQs extraídas pela plataforma, construção do checkpoint e redação das mensagens de WhatsApp.

\subsection{Otimização de checkpoints e análise de custos}

Uma das atividades centrais foi a otimização de campanhas em produção visando redução do custo operacional. O procedimento padrão consistiu em agregar os custos da campanha por sub-agente (a partir das planilhas \texttt{.xlsx} exportadas pela plataforma), identificar o modelo de linguagem mais oneroso e enxugar o checkpoint removendo apenas o conteúdo já coberto pelas FAQs, sem alterar o comportamento ou os gates operacionais do agente. A leitura dos dados e o cruzamento entre custo e configuração do agente foram feitos com Python e a biblioteca \texttt{openpyxl}.

Em um caso prático (campanha SDR Lucas Firmino -- D'Leon), o checkpoint foi reduzido de 22.230 para 11.834 caracteres (cerca de 47\%), com projeção de economia de R\$ 280 a R\$ 370 por janela de 19 dias. Em outra campanha (suporte Falcão das Milhas), o checkpoint foi reduzido em aproximadamente 58\%, mantendo todos os gates de transferência para atendimento humano.

\subsection{Integrações entre sistemas via webhook}

A última frente do estágio envolveu a integração da plataforma AwSales a sistemas externos via \textit{webhooks} HTTPS e ferramentas de orquestração visual. O ponto técnico de maior impacto foi a refatoração de uma integração de geração de \textit{deep link} de acesso a uma área de membros. A versão anterior, composta por duas \textit{tools} encadeadas, propagava erros HTTP 404 da API externa como falha técnica, disparando handoff indevido para humano. A nova versão centraliza a lógica em um \textit{workflow} no n8n, que sempre devolve HTTP 200 com um \textit{payload} indicando sucesso ou falha de aplicação, transformando a condição ``membro não encontrado'' em um caminho conversacional tratado normalmente pelo checkpoint.

\subsection{Documentação técnica}

Em paralelo às atividades técnicas, foi mantido um repositório de documentação interna em arquivos Markdown versionados em Git, contendo regras universais de criação de campanhas, configuração de \textit{tools} personalizadas, estado vivo de cada campanha e padrões de formatação obrigatórios da plataforma.

% =============================================================
% 4. RESULTADOS OBTIDOS
% =============================================================
\section{Resultados obtidos}

Os resultados alcançados ao longo dos quatro meses de estágio foram satisfatórios e atendem aos objetivos previstos no Plano de Atividades. Entre eles, destacam-se:

\begin{itemize}[leftmargin=2em]
  \item Domínio operacional da plataforma AwSales, com capacidade de configurar e otimizar campanhas completas de forma autônoma.
  \item Contribuição direta em otimizações de campanhas em produção, com redução de cerca de 47\% no tamanho do checkpoint da campanha SDR Lucas Firmino e 58\% no checkpoint da campanha de suporte Falcão das Milhas, preservando o comportamento e os gates operacionais.
  \item Refatoração de integração crítica de \textit{deep link}, com substituição de duas \textit{tools} legadas por uma única integração via n8n e padronização do contrato de resposta, reduzindo handoffs indevidos para atendimento humano.
  \item Consolidação de um repositório de documentação interna versionado em Git, organizando regras universais de criação de campanhas, configuração de \textit{tools} e estado vivo de cada cliente.
\end{itemize}

Durante o estágio, foi possível ampliar significativamente o conhecimento sobre arquiteturas de sistemas baseados em LLMs, padrões de busca semântica (RAG), engenharia de \textit{prompt}, integração entre sistemas via APIs REST e \textit{webhooks} e práticas de versionamento de artefatos textuais. A oportunidade também permitiu o contato com áreas de \textit{Customer Success}, \textit{Account Management} e Tech da empresa, ampliando a rede de contatos profissionais e abrindo possibilidades futuras na AwSales LTDA.

% =============================================================
% 5. PARTICIPAÇÃO DO ESTAGIÁRIO
% =============================================================
\section{Participação do estagiário}

\subsection{Disciplinas úteis na jornada}

A fundamentação teórica adquirida ao longo da graduação demonstrou-se indispensável para o desempenho satisfatório das atividades na empresa. A disciplina de \textbf{Programação} foi pré-requisito direto para a automação de tarefas de leitura e agregação de planilhas de custo em Python, manipulação de arquivos Markdown das campanhas e validação de formato \texttt{@tool} dos \textit{checkpoints} por expressões regulares. \textbf{Sistemas Lineares} e \textbf{Controle de Sistemas} forneceram o vocabulário e a intuição necessários para modelar a campanha como um sistema em malha fechada, onde a mensagem do lead funciona como entrada, a cadeia de sub-agentes age como controlador e atuador e as métricas de desempenho realimentam o projeto. Esse arcabouço foi essencial para a escolha de quais partes do \textit{checkpoint} eram efetivamente "lei de controle" e quais eram "modelo da planta" e deveriam ser deslocadas para a base de FAQs.

Complementarmente, \textbf{Probabilidade e Estatística} foi aplicada na análise dos dados de custo e desempenho das campanhas, com separação de variação aleatória de tendência sistemática. \textbf{Métodos Numéricos} forneceu a base conceitual para raciocinar sobre redução de custo sob restrições de qualidade. \textbf{Redes de Computadores} deu suporte ao trabalho com APIs REST, \textit{webhooks} HTTPS e tratamento de \textit{timeouts}, idempotência e contratos de mensagem. \textbf{Bancos de Dados} subsidiou o raciocínio sobre indexação semântica de FAQs (RAG) e versionamento de artefatos. \textbf{Engenharia Econômica} foi necessária para a análise de custo por campanha e o raciocínio de \textit{trade-off} técnico-econômico na escolha dos modelos de linguagem por sub-agente. Por fim, \textbf{Instrumentação} contribuiu para a leitura crítica dos \textit{dashboards} de desempenho, dos "sensores virtuais" (logs estruturados) e das métricas operacionais coletadas em produção.

\subsection{Conhecimentos complementares necessários}

Além das disciplinas da grade curricular, a execução plena das atividades demandou a aquisição de competências técnicas complementares não totalmente cobertas pelo curso. Foi essencial o estudo da arquitetura de sistemas baseados em modelos de linguagem de grande porte (LLMs), incluindo os conceitos de \textit{Retrieval-Augmented Generation} (RAG), busca semântica sobre bases vetoriais, engenharia de \textit{prompts} e orquestração de cadeias multiagentes. Também foi necessário o aprofundamento em práticas de versionamento de artefatos textuais com Git e \textit{Conventional Commits}, no uso de ferramentas de orquestração visual como n8n, Zapier e ActiveCampaign, e na compreensão dos padrões internos de integração da plataforma AwSales (formato \texttt{@tool}, padrão \textit{payload-ok} para \textit{webhooks}, restrições de formatação de \textit{checkpoints}).

Esse aprendizado foi conduzido por meio da leitura sistemática da documentação interna da empresa, do acompanhamento direto do time técnico e da inspeção de campanhas em produção. Como sugestão de conteúdo complementar a ser oferecido pela universidade, ressalta-se a relevância de uma disciplina voltada a sistemas baseados em IA e integrações entre serviços, dado o crescimento acelerado da demanda por esse perfil no mercado.

\subsection{Contribuições e inovações}

As principais contribuições do estagiário ao longo do estágio concentraram-se em três frentes complementares. A primeira foi a execução prática das otimizações de campanha, com destaque para a redução de aproximadamente 47\% no tamanho do \textit{checkpoint} da campanha SDR Lucas Firmino (D'Leon), preservando o comportamento e os \textit{gates} operacionais, e a redução de cerca de 58\% no \textit{checkpoint} da campanha de suporte Falcão das Milhas, ambos validados empiricamente em janelas de monitoramento de sete dias.

A segunda foi a consolidação metodológica, traduzida na escrita de regras universais que passaram a ser observadas em todas as campanhas novas, incluindo o formato canônico de menção a \textit{tools} (\texttt{Utilize a tool para [ação] @nome\_da\_tool}), a separação rígida entre \textit{checkpoint} (comportamento) e FAQs (conhecimento), as restrições de formatação dos \textit{checkpoints} e as regras de abertura de janela de WhatsApp.

A terceira foi o auxílio na evolução da camada de integrações, com destaque para a refatoração da \textit{tool} de \textit{deep link} de acesso à área de membros: duas \textit{tools} legadas, encadeadas diretamente contra a API externa, foram substituídas por uma única integração via \textit{workflow} no n8n, que devolve sempre HTTP 200 com um \textit{payload} bem-formado, transformando falhas externas em caminhos conversacionais normais e eliminando \textit{handoffs} indevidos para atendimento humano.

\subsection{Atividades que mais contribuíram para o desenvolvimento acadêmico e profissional}

Entre as atividades realizadas, três se destacaram pelo impacto direto no desenvolvimento acadêmico e profissional do estagiário. A primeira foi a participação no ciclo completo de diagnóstico e otimização da campanha SDR Lucas Firmino, que exigiu agregar custos por sub-agente em planilhas \texttt{.xlsx} usando Python e \texttt{openpyxl}, cruzar com a configuração do agente para identificar o gargalo (modelo Gemini 2.5 Flash respondendo por 48\% do gasto) e enxugar o \textit{checkpoint} sem perder inteligência comportamental. Esse trabalho consolidou a percepção de que problemas de operação de IA podem (e devem) ser tratados com método de engenharia.

A segunda foi a refatoração da integração de \textit{deep link} via n8n, que exigiu raciocínio sistemático sobre tratamento de erros em sistemas distribuídos, contratos de \textit{API} e mecanismos de \textit{fallback}, em estreita analogia com conceitos de instrumentação industrial e \textit{condicionamento de sinal}.

A terceira foi a consolidação do conjunto de documentos universais que normatizam a criação de campanhas na empresa (\textit{Prompt do Sistema Universal}, \textit{Estruturas e Exemplos}, \textit{Follow-Up Inteligente}, \textit{Playbook de Relatórios de Suporte e Handoff}, \textit{Guia para Criar Tool Personalizada}). Esse trabalho desenvolveu, na prática, competências de escrita técnica, governança de conhecimento e padronização de processos, hoje cada vez mais demandadas em ambientes de engenharia que envolvem sistemas baseados em IA.

% =============================================================
% 6. CONCLUSÃO
% =============================================================
\section{Conclusão}

A experiência de estágio na AwSales LTDA foi enriquecedora e proporcionou um ambiente prático para aplicar os conhecimentos teóricos adquiridos ao longo do curso de Engenharia de Controle e Automação na Universidade Federal de Itajubá -- Campus Itabira. Sob a supervisão do Gerente de Operações, Daniel Roscoe Oliveira, foi possível integrar-se à empresa e participar ativamente de suas atividades.

Durante o estágio, desde a integração inicial até a elaboração da documentação técnica final, foram aplicados conhecimentos em programação, análise de dados, sistemas de controle, integração entre serviços e engenharia econômica. Destaca-se especialmente a importância das atividades de otimização de campanhas, em que foi possível identificar gargalos de custo, propor reduções no tamanho dos checkpoints e validar empiricamente os ganhos, garantindo eficiência operacional e qualidade de atendimento.

Ao longo do processo, foram adquiridas habilidades significativas na utilização da plataforma AwSales, em ferramentas de automação como n8n e Zapier, em linguagens auxiliares como Python e em práticas de versionamento com Git. Além disso, a oportunidade permitiu ampliar a rede de contatos profissionais e fortalecer a expertise em arquiteturas de sistemas baseados em Inteligência Artificial.

Este estágio não apenas complementou a formação acadêmica, mas também abriu portas para futuras oportunidades no mercado de trabalho. Fica registrado o agradecimento pela experiência e pelas lições aprendidas durante este período na AwSales LTDA.

% =============================================================
% 7. REFERÊNCIAS BIBLIOGRÁFICAS
% =============================================================
\section{Referências Bibliográficas}

\begin{itemize}[leftmargin=2em]
  \item AWSALES. \textbf{Site Institucional}. Disponível em: \url{https://www.awsales.io/}. Acesso em: 25 maio 2026.
  \item AWSALES. \textbf{AwSales para WhatsApp -- Meta Business Partner}. Disponível em: \url{https://awsales.io/meta}. Acesso em: 25 maio 2026.
  \item BRASIL. \textbf{Lei n.\textsuperscript{o} 11.788, de 25 de setembro de 2008}: dispõe sobre o estágio de estudantes. Brasília: Presidência da República, 2008.
  \item UNIVERSIDADE FEDERAL DE ITAJUBÁ. \textbf{Diretrizes para Realização de Estágio Acadêmico do Curso de Engenharia de Controle e Automação -- Campus Itabira}. Itabira: UNIFEI. Disponível em: \url{https://sigaa.unifei.edu.br/sigaa/verProducao?idProducao=1766561}. Acesso em: 25 maio 2026.
  \item UNIVERSIDADE FEDERAL DE ITAJUBÁ. \textbf{Modelo de Relatório Final de Estágio}. Itajubá: UNIFEI. Disponível em: \url{https://graduacaoitabira.unifei.edu.br/coordenacao-geral-de-estagios/}. Acesso em: 25 maio 2026.
\end{itemize}

\end{document}
